%% kapa: exam=1 round=20 year=2009 period=1 form=A subject=civil_law subject_ko=민법 seq=12 number=12
%% source: Q-Net 감정평가사 기출문제 (cst00309) artlSeq=5024178 file=extracted/1008/5024178/2009년 제20회 감정평가사 1차시험 기출문제/1교시.hwp
%% conversion: kordoc 4.24.51 -> tools/gen1.py (auto); review=proofread-sonnet; flags=none
\begin{question}{12}
부동산등기의 유효요건에 관한 설명으로 옳지 \uline{않은} 것은? (다툼이 있으면 판례에 의함)
\begin{choices}
  \item 소유권이전등기가 된 후에 등기부가 멸실되었더라도 물권변동의 효력이 상실되는 것은 아니다.
  \item 등기는 물권변동의 효력 발생요건일 뿐 존속요건이 아니므로, 등기가 불법 말소되었더라도 물권의 존속에는 영향이 없다.
  \item 매매계약이 무효ㆍ취소되어 등기명의를 회복하는 방법으로 말소등기 대신 새로운 이전등기를 청구할 수 있다.
  \item 무효등기의 유용은 그 등기를 유용하기로 하는 합의가 이루어지기 전에 등기상 이해관계가 있는 제3자가 생기지 않은 경우에 한하여 허용된다.
  \item 등기상 이해관계인이 없는 한 멸실된 건물의 보존등기를 멸실 후 신축한 건물의 보존등기로 유용하는 것도 허용된다.
\end{choices}
\end{question}
\answer{5}
